\section{A Brief Quantum Primer}
  Quantum theory arose as a branch of physics, and so much of the
    standard notational conventions are hand-wavy, and often inconsistently
    implemented across the literature. For the sake of clarity and
    self-containedness, we specify all our
    conventions below, providing a brief summary of the
    building blocks that make up quantum computing theory along the way;
    for more details, see the celebrated introductory text by Nielsen and
    Chuang~\cite{NieChu10}.

  \todoenum[Provide explanations of:]{
    \item Bare-bones background on quantum states and Dirac notation; states
      vs.\ labels.
    \item $N = 2^n$; the set $[n]$; $n$-length registers are described by
      vectors in an $N$-sized Hilbert space.
    \item Greek letters: When and why I use them.
    \item Precise notational conventions regarding registers.
    \item Precise conventions regarding fonts.
    \item Notation for unitary evolutions: $\apply{U}$.
    \item Stick to pure states, and
      only mention the generalization to density matrices.
    \item Basic background on game hopping and advantages.
    \item Maybe write something about asymptotic vs.\ concrete security.
  }
  
  \paragraph{Dirac notation.}
    In general, quantum states are complex-valued vectors in a Hilbert space (a
      vector space defined over functions),
      which may or may not be finite-dimensional.\footnote{
        For example, consider the state $\ket{p(x)}$
        describing the position $x \in [0,1]$ of particle $p$; this will live
        in an infinite-dimensional Hilbert space characterized by the basis vectors
        $\{\ket{p(x)}\}_{x\in[0,1]}$.
      }
    %
    Quantum states are represented using the Dirac notation, in which a quantum state
      carrying the \emph{label} $\psi$ is written $\ket{\psi}$, its conjugate
      transpose is written $\bra{\psi} \coloneqq \ket{\psi}^\dagger \coloneqq
      (\ket{\psi}^*)^T$, and its
      inner product with itself is written $\braket{\psi|\psi} \coloneqq
      \abs{\ket{\psi}}^2$. Quantum states are unit-norm, and so we have that
      $\braket{\psi|\psi} = 1$. For two not necessarily equal quantum states
      $\ket{\psi}$ and $\ket{\varphi}$, their inner product
      $\braket{\psi|\varphi}$ is a real number between zero and one, where
      ``zero'' means the states are orthogonal, and ``one'' means the states are
      identical (modulo an unobservable global phase factor $e^{i\theta}$).

    This calculus is fine, but one may wonder, what \emph{is} a quantum
      state---in the sense, what kind of mathematical object is it? The answer
      can be found in the statement that they live in $N$-dimensional \emph{Hilbert
      space:} Hilbert spaces are vector spaces defined over \emph{functions} in
      place of vectors, where the inner product between two functions $f$ and
      $g$ on the interval $[a,b]$ is defined to be $\int_{-\infty}^{\infty}
      f(x)^* g(x) dx$. This inner product defines a norm, and a notion of
      orthogonality, from which a vector space may be constructed.
      Quantum states, then, are the functions that arise as solutions to a
      differential equation describing the dynamics of the system, i.e.
      the Schrödinger equation in the non-relativistic, non-spin case; the
      Pauli equation in the case with spin (such as electrons), and the Dirac
      equation in the relativistic case (excluding gravity). They take some
      input variable corresponding to an observable, such as the position $x$,
      and returns the \emph{amplitude} corresponding to that position; squaring
      this then gives the probability that observing the particle will yield
      $x$ as a result.

    If all of this sounds terribly complicated, then you understand why
      physicists grew so fond of Dirac's notation, which neatly swept all of
      these horrible details under the rug of conceptual clarity!

    Thankfully, things become a bit easier in the discrete case: the Hilbert
      space is now finite, and so it can be expressed in terms of a finite set of basis
      ``vectors'' (i.e., functions), and the inner
      product is now a sum rather than an integral. Further, functions now have finite domains,
      and may thus be written as $N$-length vectors instead, where each element
      corresponds to an input to the function, enabling us to
      employ our toolbox from linear algebra to study them. This is the
      approach taken in the quantum information sciences.

  \paragraph{Qubit basis states.}
    We are interested in finite-dimensional quantum states, implemented
      on a collection of $n$ qubits. A \emph{qubit} is a quantum state existing
      in the two-dimensional Hilbert space characterized by the basis states
      $\ket{0}$ and $\ket{1}$. In physical implementations, these states may for
      example be mapped to the spin of an electron (``up''/``down''), or the two lowest energy
      states of an atom, or in general any two orthogonal quantum states, where
      two quantum states are orthogonal if and only if there exists a
      measurement procedure that perfectly distinguishes them.

    It is worth emphasizing that the ``$0$'' and ``$1$'' here are \emph{labels},
      and do not say anything about the physical system in question other than that
      it contains (at least) two orthogonal quantum states, which are being
      identified with the two basis states $\ket{0}$ and $\ket{1}$.

    Since two qubit basis states $\ket{b_1}$ and $\ket{b_2}$, for $b_1, b_2 \in
      \{0,1\}$, the two states must be either identical (up to the
      aforementioned phase factor) or orthogonal, we have that
    \[
      \braket{b_1|b_2} = \delta_{b_1 b_2}\, ,
    \]
    where $\delta_{ij}$ is the Kronecker delta, equal to $1$ if $i = j$, and $0$
      otherwise.
      
  \paragraph{Combined quantum states.}
    In general, two quantum systems (i.e.\ states describing 
      distinct physical systems) may be combined via the tensor product (also
      known as the \emph{outer} product): 
      \[
        \ket{\psi} \otimes \ket{\varphi} \coloneqq \ket{\psi \varphi}\, .
      \]
  
    \noindent
    If $\ket{\psi}$ describes an $N$-dimensional quantum system, and
      $\ket{\varphi}$ describes an $M$-dimensional quantum system, then the
      combined state $\ket{\psi \varphi}$ describing the combined system will be
      $NM$-dimensional. (Note how the dimensions are \emph{multiplied} rather
      than added, as one might naïvely expect.)
  
    Applying this to qubits, given a collection of $n$ qubits $\ket{b_i}$, each initialized to
      one of the two basis states (so that $b_i \in \bin$), we may construct
      the combined state 
      \[
        \ket{b_0 b_1 \ldots b_{n-1}} \coloneqq \ket{b_0} \otimes \ket{b_1}
          \otimes \ldots \otimes \ket{b_{n-1}}\, .
      \]

    \noindent
    Since each qubit is a two-dimensional quantum state, we see that the
      resulting combined state will have dimension $N \coloneqq 2^n$. The resulting Hilbert
      space will be characterized by the computational\footnote{
        As with any vector space, Hilbert spaces permit an infinitude of basis
        choices; the ``computational'' basis is simply the name of the choice
        wherein basis states are labelled by classical bits (or, in general,
        bitstrings).
      } basis states
    \[
      \{\ket{i}\}_{\forall i \in [N]} = \{ \ket{b_0 b_1 \ldots b_{n-1}}
        \}_{\forall b_0, \ldots, b_{n-1} \in \bin}\, .
    \]

    \noindent
    This exponentiality in the dimension of quantum states as the number of
      qubits increases sits at the core of
      the quantum advantage over classical computers. As an illustrative
      example, the Hilbert space corresponding to a three-qubit quantum system
      is characterized by the eight computational basis states:
    \begin{align*}
      \{ \ket{i} \}_{\forall i \in [2^3]} 
      &= \{\ket{000},
        \ket{001},
        \ket{010},
        \ket{011},
        \ket{100},
        \ket{101},
        \ket{110},
        \ket{111}\} \\
      &= \left\{ \begin{bmatrix} 1 \\ 0 \\ 0 \\ 0 \\ 0 \\ 0 \\ 0 \\ 0 \end{bmatrix},
        \begin{bmatrix} 0 \\ 1 \\ 0 \\ 0 \\ 0 \\ 0 \\ 0 \\ 0 \end{bmatrix},
        \begin{bmatrix} 0 \\ 0 \\ 1 \\ 0 \\ 0 \\ 0 \\ 0 \\ 0 \end{bmatrix},
        \begin{bmatrix} 0 \\ 0 \\ 0 \\ 1 \\ 0 \\ 0 \\ 0 \\ 0 \end{bmatrix},
        \begin{bmatrix} 0 \\ 0 \\ 0 \\ 0 \\ 1 \\ 0 \\ 0 \\ 0 \end{bmatrix},
        \begin{bmatrix} 0 \\ 0 \\ 0 \\ 0 \\ 0 \\ 1 \\ 0 \\ 0 \end{bmatrix},
        \begin{bmatrix} 0 \\ 0 \\ 0 \\ 0 \\ 0 \\ 0 \\ 1 \\ 0 \end{bmatrix},
        \begin{bmatrix} 0 \\ 0 \\ 0 \\ 0 \\ 0 \\ 0 \\ 0 \\ 1 \end{bmatrix} \right\}\, .
    \end{align*}

    \noindent
    As in the single-qubit case, we see that we have orthonormality over the
      computational basis states:
    \[ \braket{x|y} = \delta_{xy}\, . \]

  \paragraph{Qubit registers.}
    Given a collection of $n$ qubits, we will often find it useful to divide it
      into \emph{registers} of specified sizes, for example to separate
      distincts elements in an input. We will use subscripts to refer to specific
      registers, unless the register we are referring to is clear from context.

    As an example, let's say we want the first
      $n_1$ qubits store the variable $x$---call this ``Register 1''---and that
      the remaining $n_2$ qubits---``Register 2''---should store the variable
      $y$. Then the following are all equivalent ways of denoting the combined state:

    \[
      \ket{x\concat y} = \ket{x}_1 \otimes \ket{y}_2 \coloneqq \ket{x}_1 \ket{y}_2
        \coloneqq \ket{x, y}\, ,
    \]
    \noindent
    where $\concat$ denotes string concatenation.

  \paragraph{Greek and Roman letters.}
    We will use the letters $x$, $y$, $z$, \ldots to represent classical binary
      strings. Thus, a state $\ket{x}$ will always be understood to be a
      computational basis state. Whenever we are talking specifically about single-qubit
      computational basis states, we will use $b$ to represent a single bit, 
      and write $\ket{b}$. 

    General quantum states are traditionally labelled by letters taken from the
      middle of the greek alphabet, such as $\psi$, $\varphi$, $\xi$, etc.
      We will follow this tradition, or otherwise use a descriptive label
      relating to the specific context in which the state is being used. 

    Quantum \emph{amplitudes}, meanwhile, are written using letters taken from the start of the greek
      alphabet ($\alpha$, $\beta$, \ldots). 
      
  \paragraph{General qubit states.}
    A general single-qubit state may be written in terms of the computational
      basis states as follows:

    \[
      \ket{\psi} = \alpha_0 \ket{0} + \alpha_1 \ket{1}\, ,
    \]
    where the amplitudes $\alpha_i$ are complex numbers satisfying
      $\abs{\alpha_0}^2 + \abs{\alpha_1}^2 = 1$, (where the absolute value of
      a complex number $\alpha$ is defined as $\abs{\alpha} = \sqrt{\alpha^*
      \alpha}$).

    Whenever a qubit is in a state
      described by amplitudes $\alpha_0$ and $\alpha_1$ that are both non-zero, we
      say that the qubit is in a \emph{superposition} over the values ``$0$'' and
      ``$1$''. If we then \emph{observe} the qubit, the probability that the
      observed outcome $X$ is either $0$ or $1$ is then given by $\prob{X=b} =
      \abs{\alpha_b}^2 = \abs{\braket{\psi |b}}^2$, and after observation, the superposition will have
      ``collapsed'' to the basis state corresponding to the observed value,
      $\ket{b}$.\footnote{
        The discernible effect of observations on quantum states has led to a
        century of heated philosophical debate over how quantum theory should be
        interpreted: The term ``collapse'' used here is typically associated with the
        conservative ``Copenhagen interpretation'', championed by Niels Bohr, and commonly
        characterized as taking a ``shut-up-and-calculate'' attitude towards the problem.
        Another popular interpretation is the many-worlds interpretation due to Hugh Everett,
        and favoured by 
        father of quantum computing theory David Deutsch, among others. Instead
        of the quantum state ``collapsing'' to become ``classical'',
        this interpretation would say that both terms are ill-defined, and that
        the universe should instead be understood to be nothing but one big
        superposition state over all worlds that might have arisen from
        the initial conditions of the Big Bang; and that what actually happened
        when you made your observation, was that the quantum state describing the subsystem you
        were studying and the quantum
        state describing the subsystem containing \emph{you} became correlated
        via the process known as \emph{entanglement}.
        (In case you're wondering, I place my own bet on the \emph{relational}
        interpretation due to Carlo Rovelli, which may be described as ``many-worlds
        without the multiverse''.)
      }

    Likewise, a general, $n$-qubit quantum state may be written in terms of the computational
      basis states:

    \[
      \ket{\psi} = \sum_{\forall i \in [2^n]} \alpha_i \ket{i} = \sum_{i = 0}^{N-1} \alpha_i \ket{i}
        = \alpha_0 \ket{0 \dots 00} + \alpha_1 \ket{0 \dots 01} + \ldots +
        \alpha_{N-1} \ket{1 \dots 11}\, ,
    \]
    
    \noindent
    where the \emph{amplitudes} $\alpha_i$ again satisfy
      $\sum_{\forall i \in [2^n]} \abs{\alpha_i}^2 = 1$.

    When expanded over the basis states, the inner product of two quantum
      states $\ket{\psi} = \sum_{\forall i \in [N]} \alpha_i \ket{i}$ and $\ket{\varphi} =
      \sum_{\forall j \in [N]} \beta_j \ket{j}$ may be written

    \[
      \braket{\psi|\varphi} = \sum_{\forall i, j \in [N]} \alpha^*_i \beta_j \braket{i|j} =
      \sum_{\forall i, j \in [N]} \alpha^*_i \beta_j \delta_{ij} = \sum_{\forall i \in [N]} \alpha^*_i \beta_i\, .
    \]

    \noindent
    We see that we have recovered the inner product of complex-valued vectors, which
      is just the usual inner product of real-valued vectors with an extra complex
      conjugation.\footnote{
        This is because inner products on any vector space must satisfy
        conjugate symmetry, $\langle x, y \rangle = \langle y, x \rangle^*$;
        see \url{https://en.wikipedia.org/wiki/Inner_product_space}.
      }
      Thus we may apply the usual operational interpretation of the inner
      product, in which vectors are ``perpendicular'' when their inner product is $0$,
      and ``parallel'' when it is $1$. Together with the requirement that
      quantum states have unit norm, this means that an inner product of $1$ implies
      that the two quantum states are actually identical,
      (up to the aforementioned unobservable phase factor).
      %This allows for an
      %operational interpretation of the inner product as \dots
      %$\sqrt{\abs{\braket{\psi|\varphi}}} =\sqrt{\abs{\sum_i \alpha^*_i \beta_i}}$.

    When expanding quantum states in terms of basis states in the manner above,
      we will sometimes implicitly exclude terms for which $\alpha_i = 0$,
      writing instead 

    \[
      \ket{\psi} = \sum_i \alpha_{i} \ket{x_i} \coloneqq \sum_{\forall i \in \Ispace} \alpha_{i} \ket{x_i}\, , 
    \]
    
    \noindent
    where $\Ispace = [\ell]$ indexes the set $\{\ket{x}\, |\ x \in [N]\ \land\
      \braket{\psi|x} \neq 0\}$ for some $\ell \in \NN$.
      %is the set of bitstrings such that the amplitude $\alpha_i$ of the
      %corresponding basis state $\ket{x_i}$ is non-zero.

  \paragraph{Unitary operations.}
    Since we are working over finite-dimensional Hilbert spaces, quantum
      operations, i.e.\ evolutions of quantum states, may be represented as $N
      \times N$-dimensional matrices. What properties should they have?
    
    Clearly, any operation that we perform on a quantum state (\emph{except}
      observing them!) should yield another valid quantum state
      as a result. Since quantum states must be unit-norm, this means
      that any operation we perform on them must preserve the norm. 
      This fact turns out to be sufficient\footnote{
        Here is a short proof of this fact:       } to restrict the space of allowed operations to the unitary matrices,
      i.e.\ invertible square matrices $U$ satisfying $U^{-1} = U^\dagger$.

    \begin{exercise}\label{ex:unitarity}
      Prove that requiring operations to be norm-preserving is equivalent to
      restricting the set of legal operations to the set of unitary matrices.
    \end{exercise}

    \noindent
    We will use capital roman letters to refer to quantum operations, in
      reference to their matrix representation. Whenever the qubit(s) that the
      operation is being applied to is not clear from context, we will provide
      the index of the qubit(s) in the
      subscript, e.g.\ $U_i$ for single-qubit operations, and $U_{i,j}$ for
      two-qubit operations, where $i,j \in [N]$ and $i \neq j$ (or more
      generally, the \emph{registers} on which the operations are applied).

    It is often useful to refer to operations in terms of the effect they have
      on the starting state, and so we define $\ket{\psi} \apply{U}
      \ket{\varphi}$ to mean that we perform the operation $\ket{\psi}
      \longrightarrow U \ket{\psi} = \ket{\varphi}$.

    Some elementary unitary transformations, known as quantum \emph{gates}, include
      \[
        \Xgate = \begin{bmatrix} 0 & 1 \\ 1 & 0 \end{bmatrix}\, , \quad
        \Zgate = \begin{bmatrix} 1 & 0 \\ 0 & -1 \end{bmatrix}\, , \quad
        \Hgate = \frac{1}{\sqrt{2}} \begin{bmatrix} 1 & 1 \\ 1 & -1 \end{bmatrix}\, , \quad
        \CXgate = \begin{bmatrix} 1 & 0 & 0 & 0 \\ 0 & 1 & 0 & 0 \\ 0 & 0 & 0 & 1 \\ 0 & 0 & 1 & 0 \end{bmatrix}\, .
      \]

    \begin{exercise}\label{ex:truthtable}
      Write down ``truth tables'' for each of the above gates by writing down
      their effect on the (single-qubit resp.\ two-qubit) computational basis states.
      (E.g., $\ket{0} \apply{\Xgate} \ket{1}$.)
    \end{exercise}

    \noindent
    The ``$C$'' in $\CXgate$ is short for ``controlled'', where the value of
      the first qubit \emph{controls} the operation on the second qubit. Any single-qubit
      operation $U$ defines a two-qubit controlled version of the same
      operation, which we will refer to as $CU$.

    \begin{exercise}
      Starting with the two-qubit state $\ket{\psi} = \ket{00}$, apply the operations $\Hgate_0$ and
      $\CXgate_{0, 1}$. What are the possible outcomes of observing the resulting
      state, and what are their probabilities? Compare with the result of
      applying the single-qubit gates $\Hgate_0$ and $\Hgate_1$ instead.
    \end{exercise}

    \begin{exercise}\label{ex:HadamardBasis}
      Write a new truth table for $\Xgate$, $\Zgate$, and $\Hgate$, only now in
      terms of the alternative basis states $\ket{+} \coloneqq \frac{\ket{0} + \ket{1}}{\sqrt{2}}$, 
      $\ket{-} \coloneqq \frac{\ket{0} - \ket{1}}{\sqrt{2}}$ instead. Compare
      with the truth table from \autoref{ex:truthtable}. What has changed?
    \end{exercise}

    \begin{exercise}
      Give the matrix representation of a unitary operation that acts on the
      alternative basis from \autoref{ex:HadamardBasis} in the same way that
      $\CXgate$ acts on the computational basis.
    \end{exercise}
